% Proofs of the dependency extracts supplied to the independent verifier.

\begin{proof}[\Cref{lem:p031-continuous-section-pointwise-openness}]
Let $N$ be a neighborhood of $\sigma(z)$.  Continuity of $\sigma$ makes
$\sigma^{-1}(N)$ a neighborhood of $z$, and the section identity gives
$\sigma^{-1}(N)\subset q(N)$.  Thus $q(N)$ is a neighborhood of $z$.

For a neighborhood $U$ of $z$, the set $q^{-1}(U)$ is a neighborhood of
$\sigma(z)$, whence
\[
 (a\circ q)^*(\sigma(z))
   \leq \sup_{q^{-1}(U)}a\circ q
   =\sup_U a.
\]
Taking the infimum over $U$ gives one inequality in
\eqref{eq:p031-continuous-section-usc-regularization}.  Conversely, for
every neighborhood $N$ of $\sigma(z)$, the set $q(N)$ is a neighborhood of
$z$ and
\[
                  \sup_N(a\circ q)=\sup_{q(N)}a\geq a^*(z).
\]
Taking the infimum over $N$ proves the reverse inequality.
\end{proof}

\begin{proof}[\Cref{cor:p031-canonical-scalar-section-anchor}]
Poineau's Corollary~3.14 makes every point of $X$ universal.  By
Definition~3.2, its canonical lift to $X_{K'}$ is the point defined by the
multiplicative tensor norm on
$\mathcal H(x)\widehat\otimes_KK'$.  Corollary~3.7 makes the lift continuous
on every affinoid domain, and Remark~3.13 globalizes that conclusion to
analytic spaces.  The lift lies over $x$, so it is a continuous section of
$\pi_{K'/K}$.  Apply
\Cref{lem:p031-continuous-section-pointwise-openness}.
\end{proof}

\begin{proof}[\Cref{lem:p031-perfect-pullback-usc}]
For $y\in q^{-1}(x)$ and every neighborhood $U$ of $x$, surjectivity gives
\[
 \sup_{q^{-1}(U)}(a\circ q)=\sup_Ua,
\]
so $b(y)\leq a^*(x)$.  The upper-semicontinuous function $b$ attains a
maximum $M$ on the compact fiber $F_x=q^{-1}(x)$.  If $r>M$, every
$y\in F_x$ has an open neighborhood $N_y$ on which $a\circ q<r$.
Compactness gives a finite union $N$ of such neighborhoods containing
$F_x$.  Since $q$ is closed,
\[
                         U:=X\setminus q(Y\setminus N)
\]
is an open neighborhood of $x$ and $q^{-1}(U)\subseteq N$.  Surjectivity
then gives $a<r$ on $U$, so $a^*(x)\leq r$.  Letting $r\downarrow M$
proves $a^*(x)=M$; the same argument for every real $r$ covers
$M=-\infty$.
\end{proof}

\begin{proof}[\Cref{cor:p031-scalar-fiber-maximum}]
Scalar projection is continuous, closed, and surjective by
\cite[\S2.9]{GR25Classical}.  Its fiber over $x$ is canonically the compact
Berkovich spectrum
\[
       \mathcal M\bigl(\mathcal H(x)\widehat\otimes_KK'\bigr)
\]
\cite[Theorem~1.2.1]{Berk12}.  Apply
\Cref{lem:p031-perfect-pullback-usc}.
\end{proof}
